\chapter{12.8 坐标系中的轴对称}


\section*{一、选择题（共8小题）}


\questionlist

\item 如果点 $P(-2, b)$ 和点 $Q(a, -3)$ 关于 $x$ 轴对称，则 $a+b$ 的值是（\quad）

  \optionlist
  \item -1
  \item 1
  \item -5
  \item 5
  \end{enumerate}

\item 在平面直角坐标系中，点 $P(-3, 2)$ 关于直线 $y=x$ 对称点的坐标是（\quad）

  \optionlist
  \item (-3, -2)
  \item (3, 2)
  \item (2, -3)
  \item (3, -2)
  \end{enumerate}

\item 将 $\triangle ABC$ 的三个顶点的横坐标不变，纵坐标乘以 $-1$，则所得图形（\quad）

  \optionlist
  \item 与原图形关于 $x$ 轴对称
  \item 与原图形关于 $y$ 轴对称
  \item 与原图形关于原点对称
  \item 向 $y$ 轴的负方向平移了一个单位
  \end{enumerate}

\item 平面内点 $A(-1, 2)$ 和点 $B(-1, 6)$ 的对称轴是（\quad）

  \optionlist
  \item $x$ 轴
  \item $y$ 轴
  \item 直线 $y=4$
  \item 直线 $x=-1$
  \end{enumerate}



\item 如图，在 $3 \times 3$ 的正方形网格中由四个格点 $A, B, C, D$，以其中一点为原点，网格线所在直线为坐标轴，建立平面直角坐标系，使其余三个点中存在两个点关于一条坐标轴对称，则原点是（\quad）
  \par\smallskip\noindent
  \begin{minipage}[t]{0.44\linewidth}
  \vspace{0pt}
 \optionlist
  \item $A$ 点
  \item $B$ 点
  \item $C$ 点
  \item $D$ 点
  \end{enumerate}
  \end{minipage}
  \hfill
  \begin{minipage}[t]{0.2\linewidth}
  \questionimage{12.8-5.png}
  \end{minipage}


\item 如图，若 $\triangle A'B'C'$ 与 $\triangle ABC$ 关于直线 $AB$ 对称，则点 $C$ 的对称点 $C'$ 的坐标是（\quad）
  \par\smallskip\noindent
  \begin{minipage}[t]{0.44\linewidth}
  \vspace{0pt}
 \optionlist
  \item (0, -1)
  \item (0, -3)
  \item (3, 0)
  \item (2, 1)
  \end{enumerate}
  \end{minipage}
  \hfill
  \begin{minipage}[t]{0.2\linewidth}
  \questionimage{12.8-6.png}
  \end{minipage}


\item 如图，$\triangle ABC$ 在平面直角坐标系中第二象限内，顶点 $A$ 的坐标是 $(-2, 3)$，先把 $\triangle ABC$ 向右平移 4 个单位得到 $\triangle A_1B_1C_1$，再作 $\triangle A_1B_1C_1$ 关于 $x$ 轴对称图形 $\triangle A_2B_2C_2$，则顶点 $A_2$ 的坐标是（\quad）
  \par\smallskip\noindent
  \begin{minipage}[t]{0.44\linewidth}
  \vspace{0pt}
 \optionlist
  \item (-3, 2)
  \item (2, -3)
  \item (1, -2)
  \item (3, -1)
  \end{enumerate}
  \end{minipage}
  \hfill
  \begin{minipage}[t]{0.2\linewidth}
  \questionimage{12.8-7.png}
  \end{minipage}


\item 如图，若直线 $m$ 经过第二、四象限，且平分坐标轴的夹角，$Rt\triangle AOB$ 与 $Rt\triangle A'OB'$ 关于直线 $m$ 对称，已知 $A(1, 2)$，则点 $A'$ 的坐标为（\quad）
  \par\smallskip\noindent
  \begin{minipage}[t]{0.44\linewidth}
  \vspace{0pt}
 \optionlist
  \item (-1, 2)
  \item (1, -2)
  \item (-1, -2)
  \item (-2, -1)
  \end{enumerate}
  \end{minipage}
  \hfill
  \begin{minipage}[t]{0.25\linewidth}
  \questionimage{12.8-8.png}
  \end{minipage}

\end{enumerate}


\section*{二、填空题（共6小题）}

\blanklist[9]

\item 点 $P(3a+6, 3-a)$ 关于 $x$ 轴的对称点在第四象限内，则 $a$ 的取值范围为\_\_\_\_\_\_。

\item 若点 $P(a, b), Q(c, d)$ 两点关于直线 $x=2$ 对称，则 $a, c$ 间的关系是\_\_\_\_\_\_，$b, d$ 间关系是\_\_\_\_\_\_；若点 $P(a, b), Q(c, d)$ 两点关于直线 $y=-2$ 对称，则 $a, c$ 间的关系是\_\_\_\_\_\_，$b, d$ 间的关系是\_\_\_\_\_\_。

\item 已知 $P_1(a-1, 5)$ 和 $P_2(2, b-1)$ 关于 $x$ 轴对称，则 $(a+b)^{2019}$ 的值为\_\_\_\_\_\_。

\rightquestionimage[width=0.35\textwidth]{0.35\textwidth}{12.8-13.png}

\item 在平面直角坐标系中，点 $A$ 的坐标为 $(-2, 3)$，点 $B$ 的坐标为 $(-1, 6)$。若点 $C$ 与点 $A$ 关于轴对称，则点 $B$ 与点 $C$ 之间的距离为\_\_\_\_\_\_。

\item 在平面直角坐标系 $A$ 中，已知直线 $l: y=x$，作 $A_1(1, 0)$ 关于 $y=x$ 的对称点 $B_1$，将点 $B_1$ 向右水平平移 2 个单位得到点 $A_2$；再作 $A_2$ 关于 $y=x$ 的对称点 $B_2$，将点 $B_2$ 向右水平平移 2 个单位得到点 $A_3$；……按此规律，则点 $B_{2019}$ 的坐标是\_\_\_\_\_\_。

\item 平面直角坐标系中有一点 $A(1, 1)$，对点 $A$ 进行如下操作：

第一步，作点 $A$ 关于 $x$ 轴的对称点 $A_1$，延长线段 $A_1A$ 到点 $A_2$，使得 $2A_1A_2 = A_1A$；

第二步，作点 $A_2$ 关于 $y$ 轴的对称点 $A_3$，延长线段 $A_2A_3$ 到点 $A_4$，使得 $2A_3A_4 = A_2A_3$；

第三步，作点 $A_4$ 关于 $x$ 轴的对称点 $A_5$，延长线段 $A_4A_5$ 到点 $A_6$，使得 $2A_5A_6 = A_4A_5$；

……

则点 $A_2$ 的坐标为\_\_\_\_\_\_，点 $A_{2025}$ 的坐标为\_\_\_\_\_\_。

若点 $A_n$ 的坐标恰好为 $(4^s, 4^t)$（$s, t$ 均为正整数），请写出 $s$ 和 $t$ 的关系式\_\_\_\_\_\_。

\end{enumerate}

\section*{三、解答题（共4小题）}

\solutionlist[15]

\item 已知点 $A(2m+n, 2)$，$B(1, n-m)$，当 $m, n$ 分别为何值时，

(1) $A, B$ 关于 $x$ 轴对称；

(2) $A, B$ 关于 $y$ 轴对称。

\vspace{2cm}

\item 如图，在平面直角坐标系中，$A(-2, 2)$，$B(-3, -2)$。



(1) 若点 $D$ 与点 $A$ 关于 $y$ 轴对称，则点 $D$ 的坐标为\_\_\_\_\_\_。

(2) 将点 $B$ 先向右平移 5 个单位再向上平移 1 个单位得到点 $C$，则点 $C$ 的坐标为\_\_\_\_\_\_。

(3) 求 $A, B, C, D$ 组成的四边形 $ABCD$ 的面积。

\rightquestionimage[width=0.35\textwidth]{0.35\textwidth}{12.8-16.png}

\vspace{3cm}

\item 三角形 $ABC$ 为等腰直角三角形，其中 $\angle A = 90^\circ$，$BC$ 长为 6。

(1) 建立适当的直角坐标系，并写出各个顶点的坐标；

(2) 将 (1) 中各顶点的横坐标都加 2，纵坐标保持不变，与原图案相比，所得的图案有什么变化？

(3) 将 (1) 中各顶点的横坐标不变，将纵坐标都乘 $-1$，与原图案相比，所得的图案有什么变化？

(4) 将 (1) 中各顶点的横坐标都乘 $-2$，纵坐标保持不变，与原图案相比，所得的图案有什么变化？

\vspace{7cm}

\item 如图所示，$\triangle COB$ 是由 $\triangle AOB$ 经过某种变换后得到的图形，观察点 $A$ 与点 $C$ 的坐标之间的关系，解答下列问题：

% 插入图片：此处应有第18题图


(1) 若点 $M$ 的坐标为 $(x, y)$，则它的对应点 $N$ 的坐标为\_\_\_\_\_\_。

(2) 若点 $P(a, 2)$ 与点 $Q(-3, b)$ 关于 $x$ 轴对称，求代数式
\[\frac{1}{ab} + \frac{1}{(a-1)(b-1)} + \frac{1}{(a-2)(b-2)} + \cdots + \frac{1}{(a-10)(b-10)}\]
的值。

\rightquestionimage[width=0.35\textwidth]{0.35\textwidth}{12.8-18.png}


\end{enumerate}
